\section{Casos de uso reales de producción remota}
\label{sec:casos_uso_remi}

La producción remota ha dejado de ser una propuesta experimental para convertirse en el
modelo operativo de referencia en eventos de gran escala. En esta sección se recogen
cuatro casos de uso representativos que ilustran el rango de aplicaciones, desde eventos
masivos con presupuesto de televisión pública hasta despliegues académicos de
investigación.

\subsection{Chaos Communication Congress (C3VOC)}

El Chaos Communication Congress (CCC) es uno de los mayores congresos de tecnología y
seguridad informática del mundo, celebrado anualmente en Alemania con más de 17\,000
asistentes y múltiples salas en emisión simultánea. El equipo de producción de vídeo del
congreso, C3VOC, desarrolló Voctomix~\cite{voctomix_github} precisamente para cubrir
esta necesidad: un mezclador software de código abierto capaz de gestionar varias salas
en paralelo sin depender de hardware broadcast propietario.

Cada sala del congreso dispone de su propio servidor con una instancia de voctocore,
alimentado por cámaras que envían señal RAW~I420 sobre TCP/MKV y operado remotamente
por un realizador a través de voctogui. Las grabaciones finales y el stream en directo
se publican en el portal de media del congreso (\texttt{media.ccc.de}). Este caso de
uso es el origen y banco de pruebas principal del proyecto: Voctomix~2.0 hereda la
arquitectura cliente-servidor de voctocore originalmente diseñada para este entorno.

\subsection{Juegos Olímpicos de Tokio 2020 (OBS)}

Los Juegos Olímpicos de Tokio 2020, celebrados en 2021, constituyeron el mayor
despliegue de producción remota en la historia del olimpismo. Olympic Broadcasting
Services (OBS) implementó un modelo REMI a gran escala en el que la señal de las
cámaras de todas las sedes olímpicas se transportó a un centro de producción centralizado
en Tokio, reduciendo drásticamente el número de equipos y personal técnico desplazado
a cada instalación~\cite{obs_tokyo_2021}.

Este despliegue validó la viabilidad técnica del modelo REMI para eventos con decenas de
sedes simultáneas, señales de alta resolución y requisitos de fiabilidad broadcast. La
reducción de la huella de carbono asociada al transporte de equipos fue uno de los
argumentos adicionales que aceleraron la adopción del modelo remoto como estándar para
ediciones futuras.

\subsection{UEFA Champions League (Telefónica Servicios Audiovisuales)}

Telefónica Servicios Audiovisuales implementó producción REMI para la retransmisión de
partidos de la UEFA Champions League jugados en el estadio Santiago Bernabéu de
Madrid~\cite{tsa_remi_2022}. Las cámaras del estadio enviaban la señal mediante SRT
sobre Internet pública hasta el centro de producción remoto, donde los realizadores
operaban el sistema como si estuvieran físicamente presentes en el estadio.

Este caso de uso demuestra la madurez del protocolo SRT para contribución broadcast
en condiciones de red reales: el cifrado AES-256 garantiza la seguridad del contenido
deportivo durante el transporte, y la corrección ARQ recupera los paquetes perdidos sin
impacto perceptible en la latencia. La operación resultó transparente para el equipo de
realización, que mantuvo su flujo de trabajo habitual sin adaptaciones significativas.

\subsection{Proyecto CyberNEMO (Universidad Politécnica de Madrid)}

En el marco del proyecto europeo de investigación CyberNEMO de la UPM, Voctomix~2.0 fue
desplegado en un clúster Kubernetes de producción para validar la arquitectura de
telemetría en condiciones reales. El broker RabbitMQ del sistema fue integrado con las
colas del clúster (\texttt{vproduction-uc-media} y \texttt{vqprobe-uc-media}), y los
eventos de mezclador, cambios de fuente activa, activaciones del stream blanker y
métricas de rendimiento del servidor se publicaron en tiempo real durante sesiones de
producción audiovisual técnica.

Este despliegue constituyó la primera validación del sistema en un entorno de
investigación académica con requisitos de seguridad y monitorización propios de
infraestructura crítica, demostrando que la arquitectura desacoplada de Voctomix~2.0
puede integrarse con plataformas externas de monitorización sin modificar el núcleo
de procesamiento multimedia.

\begin{table}[H]
  \centering
  \caption{Casos de uso representativos de producción remota audiovisual.}
  \label{tab:casos_uso_remi}
  \small
  \begin{tabular}{|l|l|l|l|}
    \hline
    \textbf{Caso de uso} & \textbf{Escala} & \textbf{Tecnología clave} & \textbf{Aportación} \\
    \hline
    CCC / C3VOC          & Múltiples salas   & Voctomix, MKV/TCP    & Origen del proyecto \\
    Tokio 2020 (OBS)     & Decenas de sedes  & REMI cloud           & Validación olímpica \\
    UEFA CL (Telefónica) & Estadio único     & SRT + AES-256        & REMI sobre Internet \\
    CyberNEMO (UPM)      & Clúster K8s       & RabbitMQ, telemetría & Validación académica \\
    \hline
  \end{tabular}
\end{table}
